% !TeX root = ../se200_identity_regimes.tex

% =============================================================================
\section{The Nine Regimes at Work}
\label{sec:application}
% =============================================================================

The preceding sections derive a nine-regime lower-bound core.
This section shows how they appear in the running diagnostic settings.
The examples are illustrative.
They show that the surplus regimes in the lower-bound core are load-bearing:
they mark distinctions that ordinary transformations force
and that accountability systems must track.

% -----------------------------------------------------------------------------
\subsection{Multi-Jurisdiction Regulatory Reporting}
\label{subsec:app_regulatory}
% -----------------------------------------------------------------------------

A corporation is an obligation-bearer.
It persists across jurisdictional differences in corporate personhood, subsidiary responsibility,
delegated liability, and compliance theory.
It is therefore represented under $\OBL$.

The reporting statute is a rule.
If it is identified by what it requires, then it is represented under $\RULEC$.
A later re-codification that preserves the requirements is the same rule under $\RULEC$
because refinement preserves content-fixed identity.
If the statute is identified by its section structure, then it is represented under $\RULES$.
The same re-codification is then a different rule,
because refinement changes the structure-fixed identity basis.

Each filing, inspection, transaction, or enforcement action is an occurrence.
It is individuated by temporal realization and provenance.
It is therefore represented under $\OCC$.
Two jurisdictions may disagree about whether the occurrence constitutes compliance.
They must still be able to refer to the same occurrence.

The jurisdiction is a scope.
If it is identified by the set of regulated cases it covers, then it is represented under $\SCOPEE$.
Splitting the jurisdiction into districts that together cover the same cases preserves identity under $\SCOPEE$.
If it is identified by its named administrative partition, then it is represented under $\SCOPES$.
The same split changes the structure and therefore breaks identity under $\SCOPES$.

Compliance indicators, filings as documents, measurement reports, and audit records are descriptive records.
They are represented under $\REC$.
They can be compared, annotated, and transferred without the substrate asserting whether compliance actually occurred.

A regulated substance, facility, instrument, parcel, or asset is a plain referent.
If the relevant identity basis is the regulated locus, then it is represented under $\LOC$.
If the relevant identity basis is the object occupying or moving through that locus,
then it is represented under $\OBJ$.

The case exercises all nine regimes in the lower-bound core.
Two jurisdictions can refer to the same corporation, statute, filing, jurisdiction, record,
and facility while evaluating compliance incompatibly.
They can also disagree about whether decomposing a jurisdiction changed the jurisdiction.
That disagreement is the $\SCOPEE/\SCOPES$ distinction.
A kind-count that stops at ``scope'' would miss it.

% -----------------------------------------------------------------------------
\subsection{Legislative History Tracking}
\label{subsec:app_legislative}
% -----------------------------------------------------------------------------

A statute is a rule.
If tracked by normative content, it is represented under $\RULEC$.
An amendment or refinement that preserves what the statute requires preserves identity under $\RULEC$.
If tracked by text, section structure, or internal organization, the statute is represented under $\RULES$.
The same amendment or refinement may then produce a new structure-fixed rule.

The enactment, amendment, vote, administrative application, judicial decision,
or repeal event is an occurrence.
It is represented under $\OCC$.
The occurrence may enact, modify, interpret, or challenge the statute,
but it is not identical with the statute.

The legislature, agency, court, officer, or regulated party is an obligation-bearer
when represented as a party that can bear responsibility, authority, or accountability.
Such entities are represented under $\OBL$.
Their identity does not depend on any single enactment or application event.

Geographic, subject-matter, procedural, temporal, or institutional boundaries are scopes.
A statutory scope may be tracked extensionally by the cases it covers, yielding $\SCOPEE$.
It may also be tracked structurally by the named jurisdictional or procedural partition, yielding $\SCOPES$.
Redistricting, reclassification, or subdivision can preserve coverage while changing structure.

Legislative records, committee reports, interpretive memoranda, docket entries, and compliance histories are records.
They are represented under $\REC$.
They may record that someone asserted a causal or normative claim without
making that claim at the substrate level.

Objects acted upon by the legal regime, such as land parcels, facilities, permits,
instruments, samples, or regulated assets, are plain referents.
They may be locus-fixed under $\LOC$ or object-fixed under $\OBJ$ depending on
what identity basis the substrate declares.

The case shows why the rule split matters.
Legal practice often needs both readings.
A doctrine may remain unchanged while the text is superseded.
A text may remain identifiable while its content is interpreted differently.
The $\RULEC/\RULES$ split prevents these two identity bases from being collapsed.

% -----------------------------------------------------------------------------
\subsection{Cross-Institutional Provenance}
\label{subsec:app_provenance}
% -----------------------------------------------------------------------------

Cross-institutional provenance separates records from the things those records describe.
A dataset entry, measurement record, transfer log, annotation record, or derived metric
is represented under $\REC$.
It may be copied, annotated, refined, and extended with provenance while remaining the same record.

The sampled object, monitored site, instrument, facility, specimen, or collection point
is a plain referent.
If the identity basis is the site, location, or collection point,
then it is represented under $\LOC$.
A fork over records concerning the same site preserves the locus-fixed referent.
If the identity basis is the specimen, instrument, or physical object,
then it is represented under $\OBJ$.
A fork that produces distinct object continuations breaks object-fixed identity.

Producing institutions, custodians, laboratories, agencies, and receiving organizations
are obligation-bearers when represented as parties responsible for custody, transfer, use, retention, or analysis.
They are represented under $\OBL$.

Data-use policies, retention rules, consent terms, sharing agreements, and access restrictions are rules.
When tracked by what they require, they are represented under $\RULEC$.
When tracked by their text, version, or internal structure, they are represented under $\RULES$.

Transfers, annotations, re-analyses, approvals, denials, measurements, and custody changes are occurrences.
They are represented under $\OCC$.
They remain referable as the same events even when institutions disagree about their significance.

Institutional coverage conditions, consent scopes, policy domains, project boundaries,
and data-sharing jurisdictions are scopes.
They may be extension-fixed under $\SCOPEE$ or structure-fixed under $\SCOPES$.
A policy domain may cover the same records after decomposition while no longer having
the same internal structure.

The provenance chain can therefore remain stable while institutions disagree about
classification, causal interpretation, data quality, responsibility, or permissible use.
The $\LOC/\OBJ$ distinction keeps ``the same site'' and ``the same specimen''
from being conflated across a fork.
The record-versus-referent distinction keeps a record from being confused with the thing it records.

% -----------------------------------------------------------------------------
\subsection{Summary}
\label{subsec:app_summary}
% -----------------------------------------------------------------------------

The proof of the lower bound was given in Section~\ref{sec:bound}.
The running cases illustrate why the lower-bound core is not artificial.

Each case uses the three core regimes that are not split in this construction:
$\OBL$, $\OCC$, and $\REC$.
Each case also exercises the three witnessed split pairs:
$\LOC/\OBJ$, $\SCOPEE/\SCOPES$, and $\RULEC/\RULES$.
The surplus regimes correspond to identity distinctions
that ordinary record operations already force.

A jurisdiction that stays the same under decomposition and one that does not are both scopes.
A statute individuated by what it requires and a statute individuated by its text are both rules.
A site and the object occupying it may both be plain referents.
The kind names are the same.
The identity bases are not.

This is a core contribution of the paper.
Reference kinds supply carriers.
Transformations determine regimes.
Counting reference kinds gives six.
The transformation analysis forces at least nine identity regimes.
